\chapter{Convenções e conversões}
\label{ap:convencoes}

\section{Unidades}

A Tabela~\ref{tab:ap-unidades} relaciona as grandezas em unidades
naturais, usadas nas derivações, às grandezas físicas. A massa $m_g$ é
sempre a massa de dispersão; $h=2\pi\hbar$.

\begin{table}[htbp]
\centering
\caption{Conversão entre unidades naturais e grandezas físicas.}
\label{tab:ap-unidades}
\small
\begin{tabular}{lll}
\toprule
Grandeza & $c=\hbar=1$ & Unidades físicas\\
\midrule
Energia de repouso & $m_g$ & $m_gc^2$\\
Número de onda de massa & $\mu=m_g$ & $\mu=m_gc/\hbar$\\
Frequência de corte & $f_g=m_g/2\pi$ & $f_g=m_gc^2/h$\\
Relação de dispersão & $\omega^2=k^2+\mu^2$ & $\omega^2=c^2k^2+(m_gc^2/\hbar)^2$\\
Massa adimensional & $u=f_gT$ & $m_gc^2=uh/T$\\
\bottomrule
\end{tabular}
\end{table}

Numericamente, $f_g=2{,}418\times10^{14}\,\mathrm{Hz}\times
(m_gc^2/\mathrm{eV})$. Para $T=4{,}5$ anos, $h/T=2{,}91\times10^{-23}$~eV;
para $T=15$ anos, $h/T=8{,}74\times10^{-24}$~eV.

\section{Sinais e curvatura}

Assinatura $(+,-,-,-)$; onda plana $\propto e^{-i\omega t+ikz}$, com
$q_\mu=(-\omega,0,0,k)$ e $\partial_\mu\to iq_\mu$. O Riemann linearizado é
dado pela Eq.~\eqref{eq:pol-riemann}. Nessa convenção,
$R^{\rm TG}_{\alpha\beta\gamma\delta}=-R_{\alpha\beta\gamma\delta}$ em
relação ao apêndice Maple do TG. Com essa convenção, a equação do
desvio geodésico, escrita com componentes espaciais euclidianas, é
$\ddot\xi^i=c^2E_{ij}\xi^j$ com $E_{ij}=R_{0i0j}$; na convenção do TG, o
sinal se inverte. Postos, razões entre componentes e correlações
angulares não dependem desse sinal global.

No Capítulo~\ref{cap:pta}, a fase é escrita como
$e^{i\omega(t-\beta\hat\Omega\cdot\mathbf x/c)}$, conjugada da anterior.
Isso não altera nenhum resultado físico, mas a conjugação deve ser
aplicada de modo consistente às correlações cruzadas complexas.

\section{Tétradas}

O apêndice do TG usa a tétrada simétrica
\begin{equation}
 l^\mu=\frac{(1,0,0,1)}{\sqrt2},\quad
 n^\mu=\frac{(1,0,0,-1)}{\sqrt2},\quad
 m^\mu=\frac{(0,1,i,0)}{\sqrt2},\quad
 \bar m^\mu=\frac{(0,1,-i,0)}{\sqrt2},
\end{equation}
com $l\cdot n=1$ e $m\cdot\bar m=-1$ na assinatura adotada. O corpo do
TG (eqs.~5.12--5.13) usa $k=\sqrt2\,l$ e $l_{\rm TG}=n/\sqrt2$. Esse
reescalonamento é um \emph{boost} da tétrada: projeções com dois vetores
$n$ mudam por um fator dois. Com as mesmas orientações, as expressões
\eqref{eq:pol-np-tg} correspondem às contrações
$\Psi_2^{\rm TG}=R(n,l,n,l)/6$, $\Psi_3^{\rm TG}=R(n,l,n,m)/2$,
$\Psi_4^{\rm TG}=R(n,m,n,m)$ e $\Phi_{22}^{\rm TG}=R(n,\bar m,n,m)$,
com o sinal $R^{\rm TG}=-R$.

\section{Bases para a helicidade zero}

Sejam $P_{ij}=\delta_{ij}-\Omega_i\Omega_j$ e $L_{ij}=\Omega_i\Omega_j$,
com $P:P=2$, $L:L=1$ e $P:L=0$. A parte espacial observável da
helicidade zero é $Q/h_0=(P-2\chi L)/\sqrt3$. A razão entre os
coeficientes longitudinal e transversal depende da base escolhida:
\begin{center}
\small
\begin{tabular}{llll}
\toprule
Convenção & Base transversal & Base longitudinal & Razão\\
\midrule
Componentes de maré ($p_l/p_b$) & $E_{11}+E_{22}$ & $E_{33}$ & $-\chi$\\
Componentes espaciais & $P$ & $L$ & $-2\chi$\\
Bases de norma dois & $P$ & $\sqrt2L$ & $-\sqrt2\chi$\\
\bottomrule
\end{tabular}
\end{center}
Uma mudança de base reescala os coeficientes e as correlações parciais
de modo que a covariância total $\sum_{AB}h_Ah_B^*\Gamma_{AB}$ seja
invariante. Nenhuma mudança de base transforma dois setores escalares
independentes no modo vinculado de Fierz--Pauli, pois os termos
cruzados expressam a coerência de uma única amplitude.

\section{Espectros}

A covariância do campo é
$\langle\tilde h_A(f,\hat\Omega)\tilde h_{A'}^*(f',\hat\Omega')\rangle
=\frac{S_A(|f|)}{8\pi}\delta_{AA'}\delta(f-f')\delta^2(\hat\Omega,\hat\Omega')$.
Com a normalização $\mathcal N_T=3/(8\pi)$ da ORF tensorial, o espectro
unilateral dos resíduos é
$S^r_{ab}=\Gamma_{ab}S_h/(6\pi^2f^2)=\Gamma_{ab}h_c^2/(12\pi^2f^3)$, com
$h_c^2=2fS_h$ \autocite{cordes2025}. Para comparar com Liang e Trodden,
cujas polarizações circulares têm norma unitária, o coeficiente de
normalização daquele trabalho deve ser $3/(4\pi)$ \autocite{liang2021}.
